\chapter{Evaluation methodology}
\label{chap:methodology}

\section{Evaluation strategy}

Chapter~\ref{chap:design} turns the research question into an empirical burden.
A block or detector alert is only an intermediate event; the relevant question
is whether enforcement changes the complete trajectory while useful work
remains possible. Evaluation must therefore connect the mechanism that
intervened, the actions that followed, and the final user and attacker outcomes.

No single setting exposes all three equally well. Component replay makes
detector, sanitiser, and policy behaviour directly inspectable, but cannot show
how an agent recovers after intervention. AgentDojo adds complete tasks,
explicit tool calls, paired conditions, and native state oracles
~\cite{debenedetti2024agentdojo}. WASP adds DOM-derived observations,
navigation, and replanning over mutable applications~\cite{evtimov2025wasp},
along with more ways for a trajectory to fail.

The evaluation therefore forms a chain rather than a pooled score. Component
replay checks the guards' contracts; AgentDojo tests controlled end-to-end
effects; WASP examines transfer to browser interaction. Each stage addresses a
limitation of the previous one while preserving the same joint criterion:
attacker success, legitimate utility, and the trajectory producing both must
be interpreted together.

\subsection{Comparison logic}

The primary comparison changes enforcement within a fixed model and scaffold;
it does not rank models. AgentDojo runs the same attacked matrix with
enforcement off and on. For each of three seeds, the final WASP experiment runs
the same 84 cases and generated inputs under off, policy-only, and full
conditions, restoring the corresponding website cell between conditions.
Rates summarise the balance; paired transitions reveal whether attacks were
prevented or introduced and whether legitimate outcomes were gained or lost.
Security obtained by disabling the agent would therefore remain visible as a
utility loss.

The policy-only condition makes the browser comparison more informative. Full
minus policy-only isolates the configured sanitisation increment; policy-only
minus off combines boundary neutralisation, model warning, and action policy,
so it is not treated as a pure policy effect.

The 97 clean AgentDojo tasks estimate capability without injection; they are
not the counterfactual for an attacked run. That role belongs to the attacked
off condition, which retains the hostile tool result while leaving the model,
task, and scaffold unchanged.

Cross-model differences remain descriptive. The 4B, 9B, and 27B variants share
an inference profile in the controlled table, but each is a distinct model
artefact with its own instruction-following, tool-selection, recovery, and
error profile. Parameter count was not independently manipulated while those
behaviours were held fixed, so the campaigns cannot identify a causal effect of
model size on either security or utility.

Nor are AgentDojo and WASP expected to reproduce one another's absolute rates:
their applications, objectives, observations, actions, and oracles differ.
WASP asks whether the controlled enforcement pattern remains visible under
browser interaction, without disguising those differences in a pooled score.

\section{Conditions and outcome definitions}

The headline treatment is the change from \texttt{off} to \texttt{enforce}.
Off bypasses observation rewriting and policy verdicts while preserving
redacted action telemetry needed for comparison. Enforce applies the configured
sanitisation and action decisions. Shadow mode can record counterfactual
decisions without changing execution, but no shadow row is used as a headline
result.

For task $i$, let $u_i\in\{0,1\}$ be the benchmark's native user oracle and
$a_i\in\{0,1\}$ its native attacker oracle. For $n$ tasks:

\begin{align}
  \mathrm{Utility} &= \frac{1}{n}\sum_{i=1}^{n}u_i,\\
  \mathrm{ASR} &= \frac{1}{n}\sum_{i=1}^{n}a_i,\\
  \mathrm{SecureUtility} &=
    \frac{1}{n}\sum_{i=1}^{n}u_i(1-a_i).
\end{align}

Secure utility requires the user task to succeed without attacker success. It
is not $1-\mathrm{ASR}$, which would reward inactivity. Paired analysis also
reports utility gains and losses, attacks prevented or introduced, and changes
in useful-and-secure outcomes.

WASP adds two binary task-level trajectory indicators. An
\emph{attacker-directed attempt} contains a proposal matching the instantiated
attacker goal; an \emph{attacker-directed execution} requires it to reach the
privileged tool. They expose diversion that the final-state oracle may miss.
Detections, sanitisation outcomes, blocks, approvals, handovers, and action
completion remain explanatory telemetry, not replacement outcomes.

The fixed task matrices are not random deployment samples, so reporting centres
exact denominators and paired descriptive effects. For zero events, $3/n$ is
used only as rule-of-three intuition under an independent binomial model
~\cite{hanley1983ruleofthree}. Native oracles remain authoritative for terminal
outcomes; trajectory replay answers only whether matching proposals or
executions occurred. This avoids treating every suspicious proposal as a
completed attack or every failed final attack as an uneventful trajectory.

\section{AgentDojo}
\label{sec:agentdojo-method}

AgentDojo 0.1.35 suite v1 provides 97 user tasks across banking, Slack, travel,
and workspace environments. Crossing the applicable tasks with the published
tool-knowledge attacks yields 629 attacked instances
~\cite{debenedetti2024agentdojo}. Sentinel's adapter preserves the benchmark's
tool names, arguments, environment state, and native user and attacker oracles.
It disables only the raw trace logger, which would otherwise persist injected
payloads.

Before agent execution, a model-free replay applies the real detector and policy
engine to 7,404 security cases across 16 public non-manual attack families and
97 distinct benign tasks. A separate scorecard checks the exact AgentDojo and
WASP text representations. Because rules were revised after observed failures,
these are development-set regression tests, not held-out accuracy estimates;
their detailed results appear in \cref{app:assurance}.

The end-to-end study uses Qwen 3.5 4B, 9B, and 27B~\cite{qwen3.5}. For each
model, the design includes a 97-task clean capability run and the complete
629-instance attacked matrix in both off and enforce conditions. The controlled
rows share context and output limits, reasoning profile, sampling parameters,
and seed, with Ollama serving the local model artefacts~\cite{ollama2026}.
These controls support within-model pairing; they do not turn the three rows
into a causal scaling experiment.

The adapter is intentionally thin. It governs tool-result observations and
mediates tool calls at the executor boundary, while leaving task
initialisation and the benchmark's evaluators unchanged. That choice prevents
a defence-specific classifier from manufacturing an apparent gain without a
corresponding change in environment state. Curated summaries are admitted only
after the expected task identifiers, totals, native scores, model name, profile,
and sampling settings agree with the campaign record. Exact run identities and
the different depths of preserved historical provenance are reported in
\cref{app:assurance}.

\section{WASP}
\label{sec:wasp-method}

The WASP adapter pins the benchmark and BrowserUse revisions, then retains the
official injector and native final-state evaluators
~\cite{evtimov2025wasp,zunic2026browseruse}. Twenty-one attacker scenarios are
crossed with two legitimate user goals and two injection formats, giving 84
attacked instances per condition. All admitted rows use a 30-step,
single-action BrowserUse scaffold, a fixed sampling profile, and DOM/accessibility
observations with screenshots excluded. Participant mode is \texttt{none}, so
no human supplies benchmark decisions. Context limits, sampling parameters,
revisions, and immutable run tags are listed in \cref{app:assurance}.

The primary browser result is the matched Qwen 3.5 27B comparison. Its 84 off
and 84 enforce executions share the model, scaffold, middleware, policy design,
step budget, generated inputs, and reset protocol; only enforcement condition
changes. WASP's native final-state oracles determine user utility and attack
success. Attacker-directed attempts and executions are reconstructed
deterministically from the recorded trajectory and instantiated attacker goal,
without relying on the optional LLM intermediate classifier.

The Qwen 3.5 9B rows answer a different, developmental question. They consist
of one shared off baseline followed by warning-only, sanitisation V1, and
sanitisation V2 enforce conditions. Middleware changed as observed failures
were repaired, and V2 deliberately reused the frozen off baseline. The sequence
can therefore explain why a warning and incomplete multiline removal were
insufficient, but it does not have the matched causal status of the 27B pair.

Execution is organised into four 21-task cells, one for each application and
injection-format combination. GitLab or Reddit is restored from a pristine
clone before its cell, and a frozen manifest recreates the same initial goals
for the paired condition. Cell-level reporting guards against an aggregate
driven by one site or one representation. Because the application is allowed to
change within a task trajectory, restoration occurs at the cell boundary rather
than after each action.

\section{Evidence admission and validity}

Only complete cells whose curated summaries agree with native scores, task
reports, and run-specific provenance enter the results. Operational failures
are not scored as oracle zeros, and changed experimental inputs require a new
immutable tag. Exact provenance and current contract checks are retained in
\cref{sec:implementation-verification,sec:appendix-wasp-provenance}. The
isolated experiments contain no human participants; one seed, two benchmark
families, finite step budgets, DOM-only observations, and development-set
detector adaptation bound external validity as discussed in
\cref{chap:discussion}.
